\section{Historical and future processing}\label{sec:periods}

\subsection{Change factors}\label{sec:changefactors}
The CMIP6 files are 20-year monthly \emph{climatologies}, but exposure is the mean over years of a non-linear
quantity (a threshold exceedance). Applying a threshold to a 20-year mean is not the same as averaging yearly
exceedances, so the main calculation uses change factors (the delta method). For every cell $i$ and calendar
month $m$, the model's change relative to the WorldClim 1970--2000 normals is applied to each observed season-year
$t\in1971..2000$:
\begin{align}
T^{f}_{\max/\min,i,t,m} &= T_{\max/\min,i,t,m} + \left(\overline T^{\,\text{CMIP6}}_{\max/\min,i,m}-\overline T^{\,\text{norm}}_{\max/\min,i,m}\right), \label{eq:delta}\\
P^{f}_{i,t,m} &= P_{i,t,m}\cdot\min\!\left(3,\ \frac{\overline P^{\,\text{CMIP6}}_{i,m}}{\overline P^{\,\text{norm}}_{i,m}}\right). \label{eq:deltap}
\end{align}
Temperature changes are additive and precipitation changes multiplicative. The ratio is capped at 3 and set to 1
where the normal is below 1\,mm.

The perturbed years are processed by exactly the same code as the observed years
(Eqs.~\ref{eq:heatcell}--\ref{eq:levels}). The \emph{historical comparator} is the same 1971--2000 years without
change factors, run through the same code. It reproduces the main pipeline's 1971--2000 country series to
\histReproMaxDiff\ (largest absolute difference), so the future and historical numbers differ only through the
climate change signal. As a check, a \emph{climatology} calculation applies the thresholds directly to the 20-year
means.

\subsection{One fixed normalisation reference}\label{sec:fixedref}
All periods are divided by the same maxima: the 2001--2024 maxima of $L_{\text{heat}}$ and $L_{\text{water}}$ and
the maximum importance of the main run (\file{index\_reference.csv}).
\begin{equation}\label{eq:fixedref}
E^{\text{period}}_{\text{heat},kc} = \frac{L^{\text{period}}_{\text{heat},kc}}{\max_{k'c'} L^{2001\text{--}24}_{\text{heat},k'c'}} .
\end{equation}
Future values above the reference are not clipped, so a future exposure can exceed one. If each period were
normalised separately, a uniform increase in stress would cancel out. The fixed reference makes the scores of
different periods comparable.

The change between periods is reported as future minus historical score. It splits exactly into a heat part and
a water part, because $S$ and $\tilde I$ are held fixed:
\begin{equation}\label{eq:decomp}
\Delta\Fd = \underbrace{\tfrac12 S_T\,\Delta E_{\text{heat}}\,\tilde I}_{\text{heat part}} + \underbrace{\tfrac12 S_P\,\Delta E_{\text{water}}\,\tilde I}_{\text{water part}} .
\end{equation}
A percentage change is reported only where the historical score is at least 1\% of the largest historical pair
score, to avoid near-zero denominators.

\subsection{What is held fixed}
\begin{itemize}
\item crop sensitivity (literature values);
\item importance (2021--2023 diets);
\item crop area, irrigation and the rainfed share (SPAM~2010);
\item calendars, flowering windows and country panels.
\end{itemize}
There is no CO$_2$ fertilisation, no CO$_2$ effect on nutrient content, no adaptation and no change in diets.
The future result therefore isolates one question: how today's crops, grown where and when they are grown today
and eaten as they are eaten today, would be exposed under the projected climate of one model.

\subsection{Conditional future yield and nutrient accounting}\label{sec:futureyield}
The normalised index is not a yield loss. To estimate nutrient consequences we return to physical units and pair
each stress with its response coefficient \emph{before} any normalisation. For each pair and season-year $t$,
\begin{equation}\label{eq:loss}
\ell^{\text{heat}}_t = \min\!\left(1,\ s_T\,H_t\right),\qquad
\ell^{\text{water}}_t = \min\!\left(1,\ K_y\,L_{\text{water},t}\right),\qquad
\ell_t = 1-(1-\ell^{\text{heat}}_t)(1-\ell^{\text{water}}_t),
\end{equation}
where $H_t$ and $L_{\text{water},t}$ are the national yearly heat and water levels. The losses are averaged over the
period ($\bar\ell$) and combined multiplicatively, so a yield cannot fall below zero. The climate-driven yield
change is the relative change in attainable yield between the historical and future runs:
\begin{equation}\label{eq:dy}
\Delta Y = -100\,\frac{\bar\ell^{\,\text{fut}}-\bar\ell^{\,\text{hist}}}{1-\bar\ell^{\,\text{hist}}}\quad[\%].
\end{equation}
Only the \emph{change} is used, because the absolute water loss from the $P/\ETc$ screening ratio is not a
credible yield level. Losses ($\Delta Y<0$) are propagated through the shock accounting of
Eq.~\ref{eq:shock} on the 2021--23 supply; gains are reported but not added to supply.

\paragraph{Empirical cross-check.} A second estimate uses the observed-yield regressions fitted on 1971--2000,
$\Delta Y^{\text{emp}} = b_T\,\Delta\Teff + b_W\cdot10\,\Delta\text{MR}$:
\begin{itemize}
\item $\text{MR}=P/\ETc$, rainfed-weighted;
\item a country's coefficient is used if it has the same sign and $p<0.05$ under three detrending methods
(linear, quadratic, first differences);
\item otherwise the pooled crop coefficient is used, under the same rule;
\item otherwise the coefficient is zero (not detected).
\end{itemize}

\paragraph{Status.} The two estimates answer different questions and are compared, not averaged. The future
nutrient numbers are a \textbf{conditional screening estimate}. They are conditional on the literature damage
functions, one climate model, and fixed area, calendars and diets. They are not a forecast.
